% ch5_answers.tex -- Quick check answers for Module 5, Chapter 5.
\begin{chapteranswers}{5}{Energy in Gravitational Fields}
  % Q5.1
  \item First find $r$ from Earth's centre:
        $r = \qty{6.371e6}{m} + \qty{0.410e6}{m} = \qty{6.781e6}{m}$. Then
        \[ U = -\frac{GMm}{r} = -\frac{(\num{6.67e-11})(\num{6.0e24})(95)}{\num{6.781e6}}
             = \qty{-5.607e9}{J}. \]
        Answer: \final{$U = \qty{-5.61e9}{J}$}.
  % Q5.2
  \item We choose $U = 0$ at infinity, where the planet's gravity has died away
        to zero. To move a satellite further away, work must be done against
        gravity. So $U$ increases as $r$ increases. $U$ keeps increasing until it
        reaches 0, but it reaches 0 only at infinity. So at every real distance
        $U$ is less than 0: it is negative.
        \begin{markingcriteria}
          \criterion{States that $U = 0$ at infinity, that $U$ increases as $r$
            increases (work is done against gravity), and so $U$ is negative at
            every real distance}{3}
          \criterion{States two of these points}{2}
          \criterion{States one relevant point}{1}
        \end{markingcriteria}
  % Q5.3
  \item Use $\Delta U = GMm\left(\dfrac{1}{r_1} - \dfrac{1}{r_2}\right)$:
        \begin{align*}
          GMm &= (\num{6.67e-11})(\num{6.0e24})(500) = \qty{2.001e17}{J.m}\\
          \frac{1}{r_1} - \frac{1}{r_2} &= \frac{1}{\num{7.0e6}} - \frac{1}{\num{9.0e6}}
            = \qty{3.175e-8}{m^{-1}}\\
          \Delta U &= (\num{2.001e17})(\num{3.175e-8}) = \qty{6.35e9}{J}
        \end{align*}
        (Or: $U_1 = \qty{-2.859e10}{J}$, $U_2 = \qty{-2.223e10}{J}$, and
        $U_2 - U_1 = \qty{6.35e9}{J}$.)
        Answers: \final{$\Delta U = +\qty{6.35e9}{J}$} and
        \final{$W_\text{by gravity} = -\Delta U = \qty{-6.35e9}{J}$}. Gravity
        does negative work because the satellite moves away.
  % Q5.4
  \item $U = -\dfrac{GMm}{r}$, so $U$ is proportional to $-\dfrac{1}{r}$. When $r$
        doubles, $U$ becomes half as big: for example \qty{-8}{GJ} becomes
        \qty{-4}{GJ} (1 GJ = \qty{1e9}{J}). Half of a negative number is closer
        to zero, so \final{$U$ increases} (it becomes less negative).
  % Q5.5
  \item $r = \qty{6.371e6}{m} + \qty{2.000e6}{m} = \qty{8.371e6}{m}$ and
        $GMm = (\num{6.67e-11})(\num{6.0e24})(300) = \qty{1.2006e17}{J.m}$.
        \begin{align*}
          K &= \frac{GMm}{2r} = \frac{\num{1.2006e17}}{2 \times \num{8.371e6}} = \qty{7.171e9}{J}\\
          U &= -\frac{GMm}{r} = -\frac{\num{1.2006e17}}{\num{8.371e6}} = \qty{-1.434e10}{J}\\
          E &= U + K = \qty{-1.434e10}{J} + \qty{7.171e9}{J} = \qty{-7.17e9}{J}
        \end{align*}
        Answers: \final{$K = \qty{7.17e9}{J}$}, \final{$U = \qty{-1.43e10}{J}$},
        \final{$E = \qty{-7.17e9}{J}$}. Check: $E = -K$ and $U = 2E$ \cmark.
  % Q5.6
  \item Here the Sun is the central mass: $M = \qty{2.0e30}{kg}$, and Earth is
        the orbiting mass: $m = \qty{6.0e24}{kg}$.
        \begin{align*}
          U &= -\frac{GMm}{r} = -\frac{(\num{6.67e-11})(\num{2.0e30})(\num{6.0e24})}{\num{1.5e11}}
             = \qty{-5.336e33}{J}\\
          E &= -\frac{GMm}{2r} = \frac{U}{2} = \frac{\qty{-5.336e33}{J}}{2} = \qty{-2.668e33}{J}
        \end{align*}
        Answers: \final{$U = \qty{-5.3e33}{J}$} and \final{$E = \qty{-2.7e33}{J}$}
        (2 significant figures, like the data). $E$ is negative: Earth is
        bound to the Sun.
  % Q5.7
  \item In a circular orbit $v^2 = \dfrac{GM}{r}$, so
        $E = -\dfrac{GMm}{2r} = -\tfrac{1}{2}m\left(\dfrac{GM}{r}\right) = -\tfrac{1}{2}mv^2 = -K$.
        You do not need $M$ or $r$.
        $E = -\tfrac{1}{2}(\num{4.2e5})(\num{6.8e3})^2 = \qty{-9.710e12}{J}$.
        Answer: \final{$E = \qty{-9.7e12}{J}$} (2 significant figures, like the data).
  % Q5.8
  \item Look at where $r$ sits in each formula.
        From $v = \sqrt{\dfrac{GM}{r}}$: $v$ \textbf{decreases}, to
        $\dfrac{1}{\sqrt{2}} = 0.71$ times its old value.
        From $K = \dfrac{GMm}{2r}$: $K$ \textbf{decreases}, to $\tfrac{1}{2}$ of
        its old value.
        From $U = -\dfrac{GMm}{r}$: $U$ becomes $\tfrac{1}{2}$ of its old value
        (for example \qty{-8}{GJ} becomes \qty{-4}{GJ}). That is closer to zero,
        so $U$ \textbf{increases}.
        $E = -\dfrac{GMm}{2r}$ does the same: it becomes $\tfrac{1}{2}$ of its
        old value, closer to zero, so $E$ \textbf{increases}.
        (The period increases to $2\sqrt{2} = 2.8$ times, by Kepler's third law.)
  % Q5.9
  \item Work is $W = Fs\cos\theta$, where $\theta$ is the angle between the force
        and the motion. In a circular orbit, gravity points to the centre. This
        is at \ang{90} to the velocity. So $\cos\ang{90} = 0$, and gravity does
        \textbf{no work}. $K$ and the speed stay the same.
        In an elliptical orbit, $r$ changes. On the way in, part of gravity
        points along the motion. So gravity does positive work, and the
        satellite speeds up ($U$ falls, $K$ rises). On the way out, part of
        gravity points against the motion. So gravity does negative work, and
        the satellite slows down. Only at perigee and apogee is gravity at
        \ang{90} to the motion. The total energy $E = K + U$ stays the same all
        the way round.
  % Q5.10
  \item Use $v_\mathrm{esc} = \sqrt{\dfrac{2GM}{r}}$:
        \begin{align*}
          v_\mathrm{esc} &= \sqrt{\frac{2(\num{6.67e-11})(\num{6.42e23})}{\num{3.39e6}}}
            = \sqrt{\qty{2.526e7}{m^2.s^{-2}}} = \qty{5026}{m.s^{-1}}
        \end{align*}
        Answer: \final{$v_\mathrm{esc} = \qty{5.03e3}{m.s^{-1}}$}.
  % Q5.11
  \item \textbf{Same speed.} To escape, $\tfrac{1}{2}mv^2 = \dfrac{GMm}{r}$. The
        mass $m$ is on both sides and cancels, so both need
        $v_\mathrm{esc} = \qty{1.12e4}{m.s^{-1}}$.
        \textbf{Different kinetic energy.} The energy needed is
        $K = \dfrac{GMm}{r_\mathrm{E}}$, which depends on $m$:
        \[ K = \frac{(\num{6.67e-11})(\num{6.0e24})(10)}{\num{6.371e6}} = \qty{6.28e8}{J} \]
        for the \qty{10}{kg} probe, and \qty{6.28e10}{J} (100 times more) for the
        \qty{1000}{kg} probe.
  % Q5.12
  \item Escape velocity is for a projectile, which gets all its energy at the
        start. A rocket's engines keep doing work on it as it climbs, so its
        total energy $K + U$ keeps rising. Also, $v_\mathrm{esc} = \sqrt{2GM/r}$
        gets smaller as $r$ gets bigger. The payload escapes once its total
        energy $K + U$ reaches zero, whatever its speed is at that moment. It
        still needs at least the same energy. It just gets it gradually.
        \begin{markingcriteria}
          \criterion{Explains that escape velocity applies to a projectile, that
            the engines keep adding energy, and that the payload escapes once
            $K + U$ reaches zero (or that $v_\mathrm{esc}$ is smaller higher up)}{3}
          \criterion{States that the engines keep adding energy, and links this
            to escape}{2}
          \criterion{States one relevant point}{1}
        \end{markingcriteria}
  % Q5.13
  \item The smallest kinetic energy needed is $K = -U = \dfrac{GMm}{r}$:
        \[ K = \frac{(\num{6.67e-11})(\num{7.35e22})(40)}{\num{1.737e6}} = \qty{1.129e8}{J}. \]
        Check: $\tfrac{1}{2}(40)(2376)^2 = \qty{1.13e8}{J}$, using the escape
        velocity from \cref{ch5:we:moon} \cmark.
        Answer: \final{$K = \qty{1.13e8}{J}$}.
  % Q5.14
  \item Use $\Delta E = \dfrac{GMm}{2}\left(\dfrac{1}{r_1} - \dfrac{1}{r_2}\right)$:
        \begin{align*}
          \frac{GMm}{2} &= \frac{(\num{6.67e-11})(\num{6.0e24})(500)}{2} = \qty{1.0005e17}{J.m}\\
          \frac{1}{r_1} - \frac{1}{r_2} &= \frac{1}{\num{8.0e6}} - \frac{1}{\num{2.4e7}} = \qty{8.333e-8}{m^{-1}}\\
          \Delta E &= (\num{1.0005e17})(\num{8.333e-8}) = \qty{8.338e9}{J}
        \end{align*}
        Answers: \final{$\Delta E = \qty{8.34e9}{J}$},
        \final{$\Delta K = -\Delta E = \qty{-8.34e9}{J}$},
        \final{$\Delta U = 2\Delta E = +\qty{1.67e10}{J}$}.
  % Q5.15
  \item $r = \qty{6.371e6}{m} + \qty{0.500e6}{m} = \qty{6.871e6}{m}$ and
        $GMm = (\num{6.67e-11})(\num{6.0e24})(650) = \qty{2.601e17}{J.m}$.
        \begin{align*}
          E_\text{start} &= U = -\frac{GMm}{r_\mathrm{E}} = -\frac{\num{2.601e17}}{\num{6.371e6}}
              = \qty{-4.083e10}{J} && \text{at rest: } K = 0\\
          E_\text{orbit} &= -\frac{GMm}{2r} = -\frac{\num{2.601e17}}{2 \times \num{6.871e6}}
              = \qty{-1.893e10}{J}\\
          \Delta E &= \qty{-1.893e10}{J} - (\qty{-4.083e10}{J}) = \qty{2.190e10}{J}
        \end{align*}
        Answer: \final{$\Delta E = \qty{2.19e10}{J}$}. (Using $\Delta U$ alone gives
        only \qty{2.97e9}{J}. That forgets the satellite's kinetic energy in orbit,
        \qty{1.89e10}{J}.)
  % Q5.16
  \item Describe three stages.
        (1) \textbf{On the pad:} the energy is stored as chemical energy in the
        fuel. $K = 0$ and $U = -\dfrac{GMm}{r_\mathrm{E}}$.
        (2) \textbf{While the engines burn:} chemical energy changes into the
        rocket's kinetic energy and potential energy (both rise). Some also
        becomes heat, sound and the kinetic energy of the exhaust gases.
        (3) \textbf{After the engines stop:} the rocket coasts upwards and slows
        down. Its $K$ changes into $U$. Ignoring air, $K + U$ stays constant.
        \begin{markingcriteria}
          \criterion{Describes all three stages: chemical energy in the fuel;
            chemical energy to $K$, $U$ and heat/sound/exhaust while burning;
            $K$ to $U$ with $K + U$ constant while coasting}{3}
          \criterion{Describes two stages correctly}{2}
          \criterion{States one correct energy change}{1}
        \end{markingcriteria}
  % Q5.17
  \item Each forward burn adds kinetic energy for a moment. But between the burns
        the satellite coasts away from Earth, and gravity slows it down, turning
        kinetic energy into potential energy. In the new, higher circular orbit
        $v = \sqrt{GM/r}$ is smaller. In energy terms, the engines raise $E$ by
        $\Delta E$, but $U$ rises by $2\Delta E$, so $K$ must fall by $\Delta E$.
        \begin{markingcriteria}
          \criterion{Explains that the burns add $K$, that coasting outwards
            turns more $K$ into $U$, and that $v = \sqrt{GM/r}$ (or $K = -E$)
            is smaller in the higher orbit}{3}
          \criterion{Explains that gravity slows the satellite as it moves
            outwards, or that $v = \sqrt{GM/r}$ is smaller for larger $r$}{2}
          \criterion{States one relevant point}{1}
        \end{markingcriteria}
  % Q5.18
  \item \textbf{Judgement:} the claim is incorrect. A decaying satellite ends up
        moving \emph{faster}.
        Drag does slow the satellite for a moment. But drag removes mechanical
        energy (it becomes heat), so $E$ falls. Because $E = -\dfrac{GMm}{2r}$,
        $r$ falls too: the satellite drops to a lower orbit. As it drops, gravity
        does positive work on it, and it speeds up by more than drag slowed it.
        In energy terms, $K = -E$, so when $E$ falls, $K$ rises. In the lower
        orbit, $v = \sqrt{GM/r}$ is bigger.
        \begin{markingcriteria}
          \criterion{Gives the judgement (incorrect), and explains that drag
            lowers $E$ and $r$, and that $K = -E$ (or $v = \sqrt{GM/r}$) so the
            satellite speeds up}{3}
          \criterion{States that the satellite speeds up and gives one
            correct reason}{2}
          \criterion{States one relevant point}{1}
        \end{markingcriteria}
\end{chapteranswers}
